% ---------------------------------------------------------------------------
% Section: The held rough error and the closure beyond the lazy regime
%
% Drop-in section for the paper. Conventions of Part III: Pi = I, c_p = 1,
% gradient flow, E = \hat Y - Y. Cross-references to existing results are
% hard-coded by number (Corollary 1, Theorem 4, Proposition 15, ...); replace
% them with \ref{...} to your labels. The theorem environments used are
% proposition, theorem, corollary, definition, assumption, remark, proof.
% ---------------------------------------------------------------------------

\section{The held rough error and the closure beyond the lazy regime}
\label{sec:held-rough}

Proposition~20 says where every expansion about $\theta_0$ stops: a rough
direction is fitted only after a displacement of order the analyticity radius
$R$, so no series in $\theta-\theta_0$ describes the stretch of training in
which the rough part of the error is removed. This section changes the
expansion variable. The one-time theory of Part~I bounds the eigenvalue of the
kernel on a rough direction by the peak feature roughness,
$\lambda_k R(v_k)^2 \le \sigma_R^2$ (Corollary~1), and Proposition~15 gives the
exact rate of $\sigma_R$ along training. Together they say that the part of
the error that a roughness function counts as rough can lose energy only
through one scalar function of time, $\sigma_R(t)$, up to a leak from the
smooth part; and the a~priori bound on the growth of $\sigma_R$ then gives a
floor on the rough part of the error that holds at every time, with no ball
and no horizon, and that decays at most as the exponential of a cubic in $t$.
Read across a nested family of seminorms, the floor says that the error at
time $t$ is at least the energy of the initial error above a cutoff mode that
moves up as $\sigma_R$ grows. The last part of the section states the one
assumption under which the same quantities give an upper bound, and the
reduced two-variable model of ignition that follows from it. Everything
before that assumption is exact.

Throughout, $R(u)=\|Pu\|$ is a projector seminorm, with $P$ a symmetric
projector on sampled functions, so that $C_R=1$ and
$\sigma_R(\theta)^2=\lambda_{\max}\big(PK(\theta)P\big)$; the cluster seminorm
and the manifold high-pass of Section~4.1 are of this form. Write
\begin{equation}
E_r = PE, \qquad E_s = (I-P)E, \qquad \sigma(t)=\sigma_R(\theta(t)),
\label{eq:split-rs}
\end{equation}
for the rough and the smooth part of the error under $R$ and for the peak
feature roughness along the path. Since $Y$ is fixed, $\dot E_r = P\dot E$,
and $\|E_r(0)\| = R(E_0) \ge R(Y)-R(\hat Y_0)$.

\subsection{The rough part loses energy only through $\sigma_R$}

\begin{proposition}[Energy of the rough part]
\label{prop:rough-energy}
Let $N_\theta$ be continuously differentiable in $\theta$ and let $\theta(t)$
be the gradient-flow path. Then, with the notation of
Eq.~\eqref{eq:split-rs}:
\begin{enumerate}
\item[(a)] Exactly,
\begin{equation}
\frac{d}{dt}\,\|E_r\|^2 = -2\,\langle E_r, K E_r\rangle - 2\,\langle E_r, K E_s\rangle,
\qquad
0 \le \langle E_r, K E_r\rangle \le \sigma^2\,\|E_r\|^2,
\qquad
|\langle E_r, K E_s\rangle| \le \sigma\,\lambda_1^{1/2}\,\|E_r\|\,\|E_s\|.
\label{eq:rough-rate}
\end{equation}
\item[(b)] For every $\varepsilon>0$ and every $t$,
\begin{equation}
\|E_r(t)\|^2 \;\ge\; \|E_r(0)\|^2\,
\exp\!\Big(-(2+\varepsilon)\!\int_0^t \sigma(s)^2\,ds\Big)
\;-\;\frac{1}{\varepsilon}\int_0^t \lambda_1(s)\,\|E_s(s)\|^2\,
\exp\!\Big(-(2+\varepsilon)\!\int_s^t \sigma(\tau)^2\,d\tau\Big)\,ds .
\label{eq:rough-floor-path}
\end{equation}
\item[(c)] The rate at which the rough part can lose energy is bounded by the
peak feature roughness alone once the smooth part is gone:
$-\frac{d}{dt}\|E_r\|^2 \le 2\sigma^2\|E_r\|^2 + 2\sigma\lambda_1^{1/2}\|E_r\|\,\|E_s\|$.
\end{enumerate}
\end{proposition}

\begin{proof}
(a) $\frac{d}{dt}\frac12\|E_r\|^2 = \langle PE, P\dot E\rangle = \langle E_r,\dot E\rangle = -\langle E_r, KE\rangle$
by Theorem~7(a), and $E = E_r+E_s$. $K$ is positive semidefinite, so
$\langle E_r,KE_r\rangle\ge0$. For a unit $\psi$ in the range of $P$,
$R(\psi)=\|\psi\|=1$, so $\psi\in R^\circ$ and $\|\psi\|_K\le\sigma_R$ by the
definition of $\sigma_R$ (proof of Theorem~1(a)); hence
$\langle E_r,KE_r\rangle=\|E_r\|_K^2\le\sigma^2\|E_r\|^2$. For the cross
term, $\langle a,Kb\rangle = g_a\cdot g_b$, so
$|\langle E_r,KE_s\rangle|\le \|E_r\|_K\|E_s\|_K\le \sigma\|E_r\|\cdot\lambda_1^{1/2}\|E_s\|$.
(b) By (a) and $2ab\le \varepsilon a^2+b^2/\varepsilon$ with
$a=\sigma\|E_r\|$ and $b=\lambda_1^{1/2}\|E_s\|$,
\[
\frac{d}{dt}\|E_r\|^2 \;\ge\; -(2+\varepsilon)\,\sigma^2\,\|E_r\|^2 \;-\;\frac{\lambda_1}{\varepsilon}\,\|E_s\|^2 ,
\]
a linear differential inequality for the nonnegative function $\|E_r\|^2$;
Gr\"onwall's inequality in its integrating-factor form gives
Eq.~\eqref{eq:rough-floor-path}. (c) is (a) with the two bounds.
\end{proof}

Part (a) is the energy identity of Theorem~7(a) projected on the rough
subspace. Its content is that the kernel enters the energy of the rough part
only through two numbers at each instant: its Rayleigh quotient on the rough
subspace, capped by $\sigma^2$, and its coupling of the rough subspace to the
smooth one, capped by $\sigma\lambda_1^{1/2}$. Part (b) is the dynamic form of
the floor of Theorem~4. Where Theorem~8(b) charges the path linearly,
$\int\sigma_R|\dot\theta|$, Eq.~\eqref{eq:rough-floor-path} charges it
exponentially, $\exp(-(2+\varepsilon)\int\sigma^2)$, and the two are different
readings of the same fact, that the output can acquire roughness only at rate
$\sigma_R$. The second term is the leak: the smooth part of the error drives
the rough part through the off-diagonal block of the kernel, and it is small
once the smooth part has been removed, which the floor of Part~I says happens
while the rough part stays. With $\varepsilon\to0$ the exponent tends to the
rate $2\sigma^2$ of Corollary~1 and the leak coefficient grows; for a bound
that is to be read along a run, $\varepsilon$ is chosen to balance the two
terms. Along a run, $\sigma(t)$ is a top eigenvalue of $PKP$ at each saved
state and $\|E_s(t)\|$ is read off, so Eq.~\eqref{eq:rough-floor-path} is a
check with no free constant.

\subsection{The a priori floor on the rough part}

Proposition~15(b) bounds the growth of $\sigma_R$ along gradient flow by
$\Theta_R\le \sigma_R\lambda_1^{1/2}H_2\|E\|$ (with $C_R=1$). Inserted into
Proposition~\ref{prop:rough-energy}, this gives a floor on the rough part that
is computed from the initial state and holds at every time.

\begin{theorem}[The rough part is held: a priori form]
\label{thm:rough-apriori}
Let $N_\theta$ be twice continuously differentiable in $\theta$, let
$\theta(t)$ be the gradient-flow path, and let $\bar\lambda_1$ and $\bar H_2$
be the largest values of $\lambda_1(\theta(s))$ and $H_2(\theta(s))$ for
$s\le t$ (or over the ball of radius $\ell(t)$ about $\theta_0$, which
contains the path). Let $\sigma_0=\sigma_R(\theta_0)$ and
\begin{equation}
a = \bar\lambda_1^{1/2}\,\bar H_2\,\|E_0\|, \qquad
\Sigma(s) = \sigma_0 + a\,s, \qquad
I(t) = \int_0^t \Sigma(s)^2\,ds = \sigma_0^2\,t + \sigma_0\,a\,t^2 + \tfrac13\,a^2\,t^3 .
\label{eq:Sigma-I}
\end{equation}
Then $\sigma(s)\le\Sigma(s)$ for all $s\le t$, and for every $\varepsilon>0$,
\begin{equation}
\|E_r(t)\|^2 \;\ge\; \|E_r(0)\|^2\, e^{-(2+\varepsilon)I(t)} \;-\; \frac{\bar\lambda_1}{\varepsilon}\, J_s(t),
\qquad
J_s(t) = \int_0^t \|E_s(s)\|^2\,ds .
\label{eq:rough-floor-apriori}
\end{equation}
In particular $\|E(t)\|\ge\|E_r(t)\|$ obeys the same bound, and
$\|E_r(0)\|\ge R(Y)-R(\hat Y_0)$.
\end{theorem}

\begin{proof}
Let $S(t)=\sigma_0^2+2\int_0^t\Theta_R$, so that $\sigma(t)^2\le S(t)$ by
Eq.~(98) and $S'=2\Theta_R\le 2\sigma\lambda_1^{1/2}H_2\|E\|\le 2S^{1/2}\bar\lambda_1^{1/2}\bar H_2\|E_0\|$,
using $\sigma\le S^{1/2}$ and $\|E\|\le\|E_0\|$ (Theorem~7(a)). Hence
$(S^{1/2})'\le a$ and $\sigma(s)\le S(s)^{1/2}\le\sigma_0+as=\Sigma(s)$.
Insert $\sigma\le\Sigma$ into Eq.~\eqref{eq:rough-floor-path}, bound
$\lambda_1(s)$ by $\bar\lambda_1$ and the inner exponential in the leak term
by one.
\end{proof}

The theorem is the counterpart, for the energy of the rough part, of the
minimum-path-length statement of Corollary~5. Three features distinguish it.
It holds at every time with no horizon and no analyticity, since nothing is
expanded. Its decay is at most exponential in a cubic: with no rough features
at the start, $\sigma_0=0$, which is the situation of Proposition~20 for a
network with a smooth activation, the exponent is $\frac{2+\varepsilon}{3}a^2t^3$,
so the rough part is held for a time of order $a^{-2/3}$. And the kernel
enters only through $\sigma_0$, $\bar\lambda_1$ and $\bar H_2$: the
initial peak feature roughness, the top eigenvalue, and the second-derivative
size, which are the three numbers through which an architecture enters the
rest of this section. The quantity $J_s(t)$ is the time integral of the energy
of the smooth part; a priori it is bounded by the band of Theorem~10 about
$e^{-K_0t}E_0$ projected on $I-P$, and along a run it is read off.

\begin{corollary}[Minimum time to remove the rough part]
\label{cor:min-time-rough}
Under the hypotheses of Theorem~\ref{thm:rough-apriori}, if $\|E_r(t)\|\le\epsilon$
with $\epsilon^2+\bar\lambda_1J_s(t)/\varepsilon<\|E_r(0)\|^2$, then
\begin{equation}
(2+\varepsilon)\,I(t) \;\ge\; \log\frac{\|E_r(0)\|^2}{\epsilon^2+\bar\lambda_1 J_s(t)/\varepsilon},
\label{eq:min-time-rough}
\end{equation}
and in particular, when $\sigma_0=0$,
\begin{equation}
t \;\ge\; \Big(\frac{3}{(2+\varepsilon)\,a^2}\,
\log\frac{\|E_r(0)\|^2}{\epsilon^2+\bar\lambda_1 J_s(t)/\varepsilon}\Big)^{1/3} .
\label{eq:cube-root}
\end{equation}
\end{corollary}

\begin{proof}
Rearrange Eq.~\eqref{eq:rough-floor-apriori}; $I$ is increasing in $t$, and
for $\sigma_0=0$, $I(t)=a^2t^3/3$.
\end{proof}

Eq.~\eqref{eq:cube-root} is the a priori ignition time of a rough direction
from a smooth start, and it depends on the architecture through
$a=\bar\lambda_1^{1/2}\bar H_2\|E_0\|$ alone: the fastest the peak feature
roughness can grow per unit time is $\bar\lambda_1^{1/2}\bar H_2$ per unit
of error, and the rough part is held until the integrated square of that
growth is of order one. Like Corollary~5 it is one-sided; unlike it, the time
scale is a cube root rather than a square, because the roughness supply
itself grows linearly while the error is held.

\subsection{The roughness spectrum and the moving cutoff}

The floor of Theorem~\ref{thm:rough-apriori} holds for every projector
seminorm. Read across a nested family, it locates where the held error sits.
Let $P_1\succeq P_2\succeq\cdots$ be a nested family of projectors, each the
projector of a seminorm $R_m(u)=\|P_mu\|$; for the manifold high-pass of
Section~4.1, $P_m$ removes the first $m$ graph-Laplacian modes, and the range
of $P_m$ shrinks as $m$ grows. Write $\sigma_m(\theta)^2=\lambda_{\max}(P_mKP_m)$,
which is nonincreasing in $m$, and
\begin{equation}
\rho_m = \|P_m E_0\|, \qquad
I_m(t) = \sigma_m(\theta_0)^2\,t+\sigma_m(\theta_0)\,a\,t^2+\tfrac13 a^2t^3 .
\label{eq:rough-spectrum}
\end{equation}
$m\mapsto\rho_m$ is the roughness spectrum of the initial error: the energy
of $E_0$ above mode $m$. It is fixed by the dataset and the initial output,
and for a target with $\hat Y_0$ small it is the roughness spectrum of the
labels.

\begin{definition}[Cutoff mode]
\label{def:cutoff}
For $t\ge0$ and $\varepsilon>0$, the cutoff mode is
\begin{equation}
m^\star(t) = \min\{\,m : (2+\varepsilon)\,I_m(t)\le\log 2\,\},
\label{eq:cutoff}
\end{equation}
the lowest mode above which, at time $t$, the a priori bound allows the rough
part to have lost at most half its energy.
\end{definition}

Since $I_m$ is nonincreasing in $m$ and increasing in $t$, the set in
Eq.~\eqref{eq:cutoff} is an up-set in $m$ and $m^\star(t)$ is nondecreasing
in $t$: the cutoff moves up through the modes as training proceeds, and it
moves even when $\sigma_m(\theta_0)=0$, through the $a^2t^3$ term, which is
feature learning.

\begin{proposition}[The error at time $t$ is at least the energy above the cutoff]
\label{prop:cutoff-floor}
Under the hypotheses of Theorem~\ref{thm:rough-apriori}, for every $t$,
\begin{equation}
\|E(t)\|^2 \;\ge\; \|P_{m^\star(t)}E(t)\|^2 \;\ge\; \tfrac12\,\rho_{m^\star(t)}^2 \;-\; \frac{\bar\lambda_1}{\varepsilon}\int_0^t\|(I-P_{m^\star(t)})E(s)\|^2\,ds .
\label{eq:cutoff-floor}
\end{equation}
\end{proposition}

\begin{proof}
Theorem~\ref{thm:rough-apriori} with $P=P_{m^\star(t)}$ and
$e^{-(2+\varepsilon)I_{m^\star}(t)}\ge\frac12$ by the definition of $m^\star$.
\end{proof}

In words: at time $t$, the error is at least the part of the initial error
that is rougher than the cutoff mode, up to the leak, and the cutoff is set
by the roughness capacity of the network, $\sigma_m(\theta_0)$ and $a$, and
by nothing else. Without feature learning ($a=0$) the cutoff is the smallest
$m$ with $\sigma_m(\theta_0)^2t\le\log2/(2+\varepsilon)$, which is the spectral
bias of the lazy regime stated in roughness rather than in frequency; with
feature learning the cutoff rises faster, and Corollary~\ref{cor:min-time-rough}
says how much faster it can rise at most. This is the quantity that the
closure below takes as its estimate of the held rough error: for a dataset
whose roughness spectrum has a step at some mode $m_Y$, the error is held near
$\rho_{m_Y}$ until $m^\star(t)$ passes $m_Y$, and the time at which it does is
Eq.~\eqref{eq:min-time-rough} for $P=P_{m_Y}$.

\subsection{The closure: efficient roughness supply}

The statements above bound the rough part from below. An upper bound, and
with it a prediction, needs one hypothesis, and it is the one that Section~14
identified as the closure that none of its assumptions provides: the growth of
the kernel on the rough directions with the sign that learning needs. The
following is its quantitative form.

\begin{assumption}[Efficient roughness supply]
\label{ass:efficient}
There is $c\in(0,1]$ such that, along the stretch $[t_0,t]$ of gradient flow,
\begin{equation}
\langle E_r, K E_r\rangle \;\ge\; c\,\sigma_R^2\,\|E_r\|^2 ,
\label{eq:efficient}
\end{equation}
that is, the Rayleigh quotient of the kernel on the current rough residual is
at least the fraction $c$ of its largest value on the rough subspace.
\end{assumption}

The assumption says that the roughness that the features carry is carried
along the residual rather than along rough directions that carry no error. It
is the positive, quantitative form of the failure of Assumption~4 on the rough
directions, and it has a mechanism: by Proposition~16, $\dot K$ is supported
on the second-order features $\nabla^2_\theta N_\theta[g_E]$ of the current
gradient, and once the smooth part has been removed, $g_E\approx g_{E_r}$, so
the kernel grows along the second-order companions of the rough residual
itself. Whether it grows along the residual's own direction, with $c$ of order
one, is the measurement: one kernel product $\langle E_r,KE_r\rangle$ and one
top eigenvalue of $PKP$ at each saved state. The assumption is a statement
about the direction of the roughness supply, not about its amount, which
remains $\sigma_R(t)$.

\begin{proposition}[Two-sided band for the rough part]
\label{prop:two-sided}
Under the hypotheses of Proposition~\ref{prop:rough-energy} and
Assumption~\ref{ass:efficient} on $[0,t]$, for every $\varepsilon\in(0,2c)$,
\begin{equation}
\|E_r(t)\|^2 \;\le\; \|E_r(0)\|^2\,e^{-(2c-\varepsilon)\int_0^t\sigma^2}
\;+\;\frac1\varepsilon\int_0^t\lambda_1(s)\,\|E_s(s)\|^2\,e^{-(2c-\varepsilon)\int_s^t\sigma^2}\,ds,
\label{eq:rough-upper}
\end{equation}
which with Eq.~\eqref{eq:rough-floor-path} encloses $\|E_r(t)\|^2$ in a band
in which the kernel enters only through $\sigma(t)$ and the leak.
\end{proposition}

\begin{proof}
From Eq.~\eqref{eq:rough-rate} and Eq.~\eqref{eq:efficient},
$\frac{d}{dt}\|E_r\|^2\le -2c\sigma^2\|E_r\|^2+2\sigma\lambda_1^{1/2}\|E_r\|\|E_s\|
\le -(2c-\varepsilon)\sigma^2\|E_r\|^2+(\lambda_1/\varepsilon)\|E_s\|^2$,
and Gr\"onwall's inequality.
\end{proof}

Under the assumption, the rough part of the error is predicted by one scalar
function of time, $\sigma_R(t)$, to within the band set by $c$ and the leak.
What remains open is $\sigma_R(t)$ itself: Proposition~15(b) gives its rate
exactly, $\frac{d}{dt}\sigma_R^2\le2\Theta_R$ with $\Theta_R$ the largest drive
over the dual ball, and bounds it by $\sigma_R\lambda_1^{1/2}H_2\|E\|$; its
actual value needs the second derivative of the network along the path, which
is where the frozen-residual reduced flow of the last part enters. The
following reduced model is what the band becomes when both the supply and its
growth are taken at their bounds.

\begin{proposition}[Reduced model of ignition]
\label{prop:reduced}
Let $r=\|E_r\|$ and suppose that on a stretch of training the smooth part has
been removed, $E\approx E_r$, that Assumption~\ref{ass:efficient} holds with
constant $c$, and that the drive attains its bound with the favourable sign,
$\Theta_R=\sigma_R\lambda_1^{1/2}H_2\|E\|$, with $\lambda_1$ and $H_2$
treated as constants $b=\lambda_1^{1/2}H_2$. Then $(r,\sigma)$ obey
\begin{equation}
\dot r = -c\,\sigma^2\,r, \qquad \dot\sigma = b\,r ,
\label{eq:reduced-ode}
\end{equation}
which have the first integral $r+\frac{c}{3b}\sigma^3=\mathrm{const}$. Hence
the peak feature roughness that the network reaches when the rough residual
$r_0$ has been removed is
\begin{equation}
\sigma_\infty = \Big(\sigma_0^3+\frac{3b}{c}\,r_0\Big)^{1/3},
\label{eq:sigma-infty}
\end{equation}
and the time at which $r$ has fallen to $r_0/2$ is
$t_{1/2}=\int_{\sigma_0}^{\sigma_{1/2}}d\sigma\big/\big(b\,r_0-\tfrac{c}{3}(\sigma^3-\sigma_0^3)\big)$
with $\sigma_{1/2}^3=\sigma_0^3+3br_0/(2c)$; for $\sigma_0=0$ this is
$t_{1/2}=\kappa_c\,(b^2r_0^2c)^{-1/3}$ with
$\kappa_c=\int_0^{(3/2)^{1/3}}du/(1-u^3/3)$, a number.
\end{proposition}

\begin{proof}
The first equation is Proposition~\ref{prop:rough-energy}(a) with the cross
term dropped and the quadratic form at its assumed value; the second is
$\frac{d}{dt}\sigma^2=2\Theta_R=2\sigma b r$. Then
$dr/d\sigma=-c\sigma^2/b$, which integrates to the first integral; set $r=0$
for $\sigma_\infty$, and $dt=d\sigma/(br)$ with $r$ from the first integral
for $t_{1/2}$; the substitution $\sigma=(3br_0/c)^{1/3}u$ gives the stated form.
\end{proof}

Eq.~\eqref{eq:sigma-infty} is a statement of what the kernel needs: to fit a
rough residual of size $r_0$, the features must become rough to the level
$\sigma_\infty$, and the network supplies roughness at rate at most $b$ per
unit of error. It is the counterpart in roughness of Proposition~20 in
displacement, and unlike that proposition it is a model rather than a bound:
its two inputs, $c$ and the realized $\Theta_R$ against its bound, are
measured. The cube-root scaling of $t_{1/2}$ with $b^2r_0^2$ is the one of
Corollary~\ref{cor:min-time-rough}, which holds without the model; the model
adds the constant and the second half of the curve, the decay of $r$ once
$\sigma$ has grown, which is exponential at the rate $c\sigma_\infty^2$ as
$r\to0$.

\subsection{Depth and the comparison of architectures}

Everything above enters an architecture through three numbers at the start of
training, $\sigma_R(\theta_0)$, $\bar\lambda_1$ and $\bar H_2$, and through
the two run quantities $c$ and $\Theta_R$. For two designs on the same
dataset, the comparison principle of Corollary~6 applies with the bounds of
this section.

\begin{corollary}[Comparing designs on the rough part]
\label{cor:depth}
Let designs $A$ and $B$ be trained by gradient flow on the same dataset, let
$P$ be the projector of a roughness function fixed by the inputs, and let
$T>0$.
\begin{enumerate}
\item[(i)] (Proved.) $\|E_r^A(T)\|^2$ is at least the right side of
Eq.~\eqref{eq:rough-floor-apriori} computed with $\sigma_0^A$, $a_A$ and
$J_s^A(T)$.
\item[(ii)] (Under Assumption~\ref{ass:efficient} for $B$ with constant $c_B$
and with $\sigma_B(t)$ read off the run, or taken from the reduced model.)
$\|E_r^B(T)\|^2$ is at most the right side of Eq.~\eqref{eq:rough-upper}.
\item[(iii)] If the bound in (ii) is below the bound in (i), then
$\|E_r^B(T)\|<\|E_r^A(T)\|$: $B$ has learned more of the rough part of the
labels by time $T$.
\end{enumerate}
\end{corollary}

The corollary separates what depth does into two effects that can be checked
one at a time. A deeper network can have larger initial peak feature
roughness $\sigma_R(\theta_0)$, because its gradient features decorrelate over
shorter input distances when the initialization is near the edge of chaos;
that lowers the held rough error after the lazy phase, since the cutoff of
Definition~\ref{def:cutoff} starts higher. And it can have larger
$\lambda_1^{1/2}H_2$, which is the rate $b$ at which it can supply roughness
per unit of error; that shortens the ignition time of
Corollary~\ref{cor:min-time-rough} and Proposition~\ref{prop:reduced}. Both
are computed before training, one top eigenvalue of $PKP$ and one nested
automatic differentiation per design, and both are expected to grow with
depth; the theory says which part of the error each one removes and in what
time, but not that either grows with depth, which is a measurement or a
statement about the initialization. Where the initialization statistics of
deep networks are known in closed form, they can supply $\sigma_R(\theta_0;L)$,
$\lambda_1(L)$ and $H_2(L)$ as functions of depth, and this is the only place
in the theory where a limit of the initialization enters; nothing along the
path is taken to a limit.

\subsection{Scope and what to measure}

\textbf{Exact under gradient flow, no hypothesis beyond smoothness.}
Proposition~\ref{prop:rough-energy}; Theorem~\ref{thm:rough-apriori} and
Corollary~\ref{cor:min-time-rough} with $\bar\lambda_1$ and $\bar H_2$ taken
along the path or over the ball of radius $\ell(t)$;
Proposition~\ref{prop:cutoff-floor}; Corollary~\ref{cor:depth}(i).
The leak term is exact and is read off a run; its a priori bound through
Theorem~10 is available while the band of that theorem is informative.
\textbf{Under Assumption~\ref{ass:efficient}.} Proposition~\ref{prop:two-sided};
Corollary~\ref{cor:depth}(ii)--(iii). \textbf{Model.}
Proposition~\ref{prop:reduced}, which takes the supply and its growth at their
bounds. The present section breaks the rule of Part~III that nothing is
approximated, in its last two parts only, and it does so with a proved band
around the approximation.

The measurements that decide whether the closure is the right one, to be added
to the list of Section~15: (13) $\sigma_R(t)$ along each run for the
manifold high-pass at several $m$, against its a priori bound $\Sigma(t)$ of
Eq.~\eqref{eq:Sigma-I} and against the realized drive $\Theta_R$;
(14) the ratio $\langle E_r,KE_r\rangle/(\sigma_R^2\|E_r\|^2)$ along each run,
which is the constant $c$ of Assumption~\ref{ass:efficient} and the quantity
that decides whether the roughness supply is spent on the residual;
(15) $\|E_r(t)\|^2$ against the two sides of the band,
Eqs.~\eqref{eq:rough-floor-path} and \eqref{eq:rough-upper}, and $J_s(t)$;
(16) the cutoff mode $m^\star(t)$ against the mode at which the measured
error's roughness spectrum actually drops, which tests
Proposition~\ref{prop:cutoff-floor} as an estimate rather than as a bound;
(17) $\sigma_R(\theta_0)$, $\lambda_1$ and $H_2$ as functions of depth at
initialization, which are the three numbers through which depth enters
Corollary~\ref{cor:depth}.
